\section{Concurrency}

\begin{frame}{What is Concurrency?}
  \begin{block}{Definition}
    \textbf{Concurrency} is the ability of multiple users or processes to access
    and modify the same database data \textbf{simultaneously}, without compromising
    data integrity or consistency.
  \end{block}

  In any real production system -- online shop, banking app, ERP -- hundreds or
  thousands of transactions execute at the same moment. This is the \emph{normal}
  operating mode, not an edge case.

  \begin{block}{Four Concurrency Problems}
    \begin{tabular}{@{}ll@{}}
      \toprule
      \textbf{Problem} & \textbf{What happens} \\
      \midrule
      Lost Update          & One transaction's write overwrites another's \\
      Dirty Read           & Reading uncommitted data that is later rolled back \\
      Non-Repeatable Read  & Same read gives different values within one transaction \\
      Phantom Read         & A repeated query returns a different set of rows \\
      \bottomrule
    \end{tabular}
  \end{block}

  \begin{block}{Goal: Serializability}
    The result of concurrent execution must be equivalent to \emph{some} serial
    (one-at-a-time) execution of the same transactions.
  \end{block}
\end{frame}

\begin{frame}[fragile]{Problem 1: Lost Update}
  \textbf{Scenario:} two users simultaneously update the price of the same track.

  \medskip
  \begin{tabular}{@{}p{3.8cm}p{3.8cm}p{3.5cm}@{}}
    \toprule
    \textbf{Transaction A} & \textbf{Transaction B} & \textbf{What happens?} \\
    \midrule
    \texttt{SELECT price} (= 9.99) & & \\
    & \texttt{SELECT price} (= 9.99) & Both see old value \\
    Compute 9.99 + 2.00 & & \\
    & Compute 9.99 $-$ 1.00 & \\
    \texttt{UPDATE price = 11.99} & & A commits first \\
    \texttt{COMMIT} & & \\
    & \texttt{UPDATE price = 8.99} & \textbf{B overwrites A!} \\
    & \texttt{COMMIT} & \textbf{A's change is LOST} \\
    \bottomrule
  \end{tabular}

  \begin{alertblock}{Result}
    Correct result should be 10.99 (both applied). Actual result: 8.99.
    A's +2.00 update is \textbf{completely lost}.
  \end{alertblock}

  \begin{block}{SQL Scenario}
\begin{lstlisting}[style=sqlstyle]
-- Both sessions run this simultaneously on TrackID = 4
BEGIN;
  SELECT price FROM Track WHERE TrackID = 4;  -- both see 9.99
  UPDATE Track SET price = price + 2.00       -- A: 11.99
  -- or:                   price = price - 1  -- B: 8.99
COMMIT;
-- B's COMMIT runs after A's => B wins, A's change is gone
\end{lstlisting}
  \end{block}
\end{frame}

\begin{frame}[fragile]{Problem 2: Dirty Read (Uncommitted Dependency)}
  \textbf{Scenario:} Transaction A reads data from Transaction B before B commits.

  \medskip
  \begin{tabular}{@{}p{3.5cm}p{4cm}p{3.2cm}@{}}
    \toprule
    \textbf{Transaction A (ATM)} & \textbf{Transaction B (Transfer)} & \textbf{Event} \\
    \midrule
    & \texttt{BEGIN} & \\
    & Balance: 300 + 500 = 800 & In memory, not committed \\
    \texttt{SELECT balance} = 800 & & Reads \emph{dirty} data! \\
    Displays \EUR{800} to user & & \\
    & \texttt{ROLLBACK} (transfer failed) & Balance reverts to 300 \\
    \multicolumn{3}{l}{\textit{ATM showed \EUR{800} but real balance is \EUR{300}.}} \\
    \bottomrule
  \end{tabular}

  \begin{alertblock}{Result}
    A read uncommitted data from a transaction that was later rolled back.
    The ATM customer may attempt to withdraw money they do not actually have.
  \end{alertblock}

  \begin{block}{Prevented by isolation level}
\begin{lstlisting}[style=sqlstyle]
-- Dirty reads are prevented at READ COMMITTED or higher
SET TRANSACTION ISOLATION LEVEL READ COMMITTED;
\end{lstlisting}
  \end{block}
\end{frame}

\begin{frame}[fragile]{Problem 3: Non-Repeatable Read}
  \textbf{Scenario:} an auditor reads three account balances while a transfer runs concurrently.\\
  Initial state: Acc\,1 = \EUR{400}, Acc\,2 = \EUR{300}, Acc\,3 = \EUR{700} (total = \EUR{1\,400}).

  \medskip
  \begin{tabular}{@{}p{3.5cm}p{4cm}p{3cm}@{}}
    \toprule
    \textbf{Auditor (A)} & \textbf{Transfer (B)} & \textbf{Event} \\
    \midrule
    Read Acc\,1 = 400 & & \\
    Read Acc\,2 = 300 & & \\
    & Acc\,3: 700 $-$ 600 = 100 & Transfer begins \\
    & Acc\,1: 400 + 600 = 1000 & \\
    & \texttt{COMMIT} & \\
    Read Acc\,3 = 100 & & Sees committed new value! \\
    Sum = 400+300+100 = \textbf{800} & & \textbf{Wrong! Should be 1400} \\
    \bottomrule
  \end{tabular}

  \begin{alertblock}{Result}
    Auditor A read accounts at different points in time and saw a snapshot that
    never actually existed -- \EUR{600} appears to have vanished.
  \end{alertblock}

  \begin{block}{Prevented by isolation level}
\begin{lstlisting}[style=sqlstyle]
-- Prevented at REPEATABLE READ or SERIALIZABLE
SET TRANSACTION ISOLATION LEVEL REPEATABLE READ;
\end{lstlisting}
  \end{block}
\end{frame}

\begin{frame}{Transaction Isolation Levels}
  SQL defines four isolation levels that trade off correctness vs.\ concurrency:

  \begin{center}
  \renewcommand{\arraystretch}{1.3}
  \small
  \begin{tabular}{@{}lcccc@{}}
    \toprule
    \textbf{Isolation Level} & \textbf{Dirty Read} & \textbf{Non-Rep. Read} & \textbf{Phantom Read} \\
    \midrule
    \texttt{READ UNCOMMITTED} & Possible & Possible & Possible \\
    \texttt{READ COMMITTED}   & \textcolor{green!60!black}{Prevented} & Possible & Possible \\
    \texttt{REPEATABLE READ}  & \textcolor{green!60!black}{Prevented} & \textcolor{green!60!black}{Prevented} & Possible \\
    \texttt{SERIALIZABLE}     & \textcolor{green!60!black}{Prevented} & \textcolor{green!60!black}{Prevented} & \textcolor{green!60!black}{Prevented} \\
    \bottomrule
  \end{tabular}
  \end{center}

  \begin{block}{MySQL Defaults}
    MySQL InnoDB default: \textbf{REPEATABLE READ} (with snapshot isolation, preventing most phantoms too).\\
    Most application frameworks use \textbf{READ COMMITTED} for better throughput.
  \end{block}
\end{frame}

\begin{frame}{Concurrency Control Strategies}
  Two fundamentally different philosophies for dealing with concurrent access:

  \medskip
  \begin{block}{Optimistic Concurrency Control (OCC)}
    \textbf{Assumption:} conflicts are rare. Let all transactions proceed freely.
    \begin{itemize}
      \item Transactions read and write without acquiring any locks.
      \item At \texttt{COMMIT} time, the DBMS checks whether a conflict occurred.
      \item If a conflict is detected, the transaction is \textbf{aborted and retried}.
      \item Pros: simple to implement; high throughput when conflicts are actually rare.
      \item Cons: wasted work on conflict; poor performance under high contention.
    \end{itemize}
  \end{block}

  \begin{block}{Pessimistic Concurrency Control (Lock-Based)}
    \textbf{Assumption:} conflicts are likely. Prevent them before they happen.
    \begin{itemize}
      \item A transaction must \textbf{acquire a lock} on a data item before accessing it.
      \item Other transactions that need the same item must \textbf{wait} for the lock to be
            released.
      \item Locks are released at transaction end (commit or rollback).
      \item Pros: conflicts never occur in the first place; correct even under high contention.
      \item Cons: complex; reduces parallelism; can cause \textbf{deadlocks}.
    \end{itemize}
  \end{block}
\end{frame}

\begin{frame}{Lock Granularity}
  A \textbf{lock} can be placed at different levels of the data hierarchy. The \emph{granularity}
  of a lock determines how much data is protected -- and how much concurrency is sacrificed.

  \medskip
  \begin{tabular}{@{}p{2cm}p{4.5cm}p{4cm}@{}}
    \toprule
    \textbf{Level} & \textbf{Description} & \textbf{Typical Use Case} \\
    \midrule
    Database & Locks the entire database; only one user can access it at a time &
               System maintenance, bulk import, schema migrations \\[4pt]
    Table    & Locks a single table; all rows are blocked for other writers &
               Few-user systems; batch operations on one table \\[4pt]
    Row (Tuple) & Locks a single row; other rows in the same table are freely accessible &
               Standard production OLTP systems (MySQL InnoDB, PostgreSQL) \\[4pt]
    Cell (Field) & Locks a single attribute value within a row &
               Extremely rare; very complex to implement \\
    \bottomrule
  \end{tabular}

  \medskip
  \begin{block}{Design principle}
    Always lock \textbf{the smallest amount of data necessary} to maintain correctness. Coarser
    locks are easier to implement but destroy concurrency. Fine-grained row-level locks maximise
    parallelism but require more overhead to track.
  \end{block}
\end{frame}

\begin{frame}{Shared and Exclusive Locks}
  \begin{block}{Exclusive Lock (X-Lock) -- Write Lock}
    Acquired by \texttt{UPDATE}, \texttt{INSERT}, and \texttt{DELETE} operations. While an X-lock
    is held on a data item, \textbf{no other transaction may acquire any lock} (neither shared nor
    exclusive) on that item. Writes require exclusive access.
  \end{block}

  \begin{block}{Shared Lock (S-Lock) -- Read Lock}
    Acquired by \texttt{SELECT} operations. Multiple transactions may hold S-locks on the same
    item simultaneously -- concurrent reads are always safe. However, \textbf{an X-lock cannot be
    granted} while any S-lock is held, because a write while others are reading would produce
    inconsistent results.
  \end{block}

  \medskip
  \textbf{Lock compatibility matrix:}

  \begin{tabular}{@{}lcc@{}}
    \toprule
    \textbf{Requested} $\backslash$ \textbf{Held} & \textbf{S-Lock held} & \textbf{X-Lock held} \\
    \midrule
    Request S-Lock & \textcolor{green!60!black}{Granted} & \textcolor{red}{Wait} \\
    Request X-Lock & \textcolor{red}{Wait} & \textcolor{red}{Wait} \\
    \bottomrule
  \end{tabular}

  \medskip
  \begin{examples}{Summary}
    Reads can run in parallel (S+S is fine). A write blocks all other accesses (X blocks
    everything). This implements the intuitive rule: many readers are fine; a writer needs
    exclusive access.
  \end{examples}
\end{frame}

\begin{frame}[fragile]{Deadlocks}
  \begin{alertblock}{Definition}
    A \textbf{deadlock} occurs when two transactions are each waiting for the
    other to release a lock -- a circular dependency neither can escape alone.
  \end{alertblock}

  \begin{block}{Classic SQL Deadlock Example}
\begin{lstlisting}[style=sqlstyle]
-- Session A                      -- Session B
BEGIN;                             BEGIN;
UPDATE Track                       UPDATE Track
  SET Duration = 300                 SET Duration = 400
  WHERE TrackID = 1;                 WHERE TrackID = 2;  -- OK
                                   UPDATE Track
UPDATE Track                         SET Duration = 500
  SET Duration = 350                 WHERE TrackID = 1;  -- WAITS for A
  WHERE TrackID = 2; -- WAITS for B
-- DEADLOCK: A waits for B, B waits for A
-- MySQL detects the cycle and kills one transaction (ERROR 1213)
\end{lstlisting}
  \end{block}

  \begin{block}{Resolution Strategies}
    \begin{itemize}
      \item \textbf{Detection}: DBMS builds a ``waits-for'' graph; when a
            cycle is found, one transaction (the \emph{victim}) is aborted
            and retried. Used by MySQL InnoDB and PostgreSQL.
      \item \textbf{Prevention}: always acquire locks in the \emph{same
            order} across all transactions. No cycle can form if all
            sessions lock TrackID 1 before TrackID 2.
    \end{itemize}
  \end{block}
\end{frame}
